\appendix
\chapter{ProVerif Formal Verification Artefacts}
\label{appx:proverif}

This appendix reproduces the ProVerif source models and the corresponding tool output transcripts that underpin the formal verification results summarised in Table~\ref{tab:proverif} of Chapter~\ref{ch:results}. The three models target ballot privacy, individual verifiability, and voter eligibility, respectively; each is verified by ProVerif~2.05~\cite{blanchet2016} under the symbolic Dolev--Yao adversary model. All three queries return \texttt{true}.

The model files are versioned in the project repository under the directory \texttt{proverif/}. The shell script \texttt{proverif/run\_verification.sh} reproduces all three runs in sequence and prints the \texttt{RESULT} lines.

\section{Reproducing the Verification}
\label{sec:appx-proverif-repro}

Install ProVerif via OPAM (OCaml package manager):
\begin{verbatim}
opam install proverif
\end{verbatim}

Run all three verifications:
\begin{verbatim}
cd proverif
./run_verification.sh
\end{verbatim}

The expected condensed output is:
\begin{verbatim}
RESULT Observational equivalence is true.
RESULT inj-event(BallotRecorded(b_2)) ==>
       inj-event(VoterCastBallot(pkV,b_2)) is true.
RESULT event(ValidBallotCast(pkV)) ==>
       event(VoterRegistered(pkV)) is true.
\end{verbatim}

\section{Ballot Privacy: \texttt{privacy.pv}}
\label{sec:appx-proverif-privacy}

Ballot privacy is verified via observational equivalence: an adversary controlling the network cannot distinguish between two honest voters swapping their ballots. The model uses the \texttt{diff[v0, v1]} construct to test whether the protocol is observationally equivalent under both ballot assignments.

\subsection*{Source: \texttt{proverif/privacy.pv}}
\begin{footnotesize}
\begin{verbatim}
(* CovertVote: Ballot Privacy via Observational Equivalence *)
(* NO queries allowed in this file -- diff mode only *)

type skey.
type pkey.
type vote.
type nonce.
type commitment.
type zkprf.
type signature.
type keyimage.

free ch: channel.
free bb: channel.
free sa2_a: channel [private].
free sa2_b: channel [private].

free v0: vote.
free v1: vote.

fun pk(skey): pkey.
fun enc(vote, pkey, nonce): bitstring.
fun commit(vote, nonce): commitment.
fun zkp_make(vote, nonce, commitment): zkprf.
fun ring_sign(bitstring, skey, pkey): signature.
fun key_image(skey): keyimage.
fun sa2_share_a(bitstring, nonce): bitstring.
fun sa2_share_b(bitstring, nonce): bitstring.
fun make_ballot(bitstring, commitment, zkprf, signature, keyimage): bitstring.

let SA2ServerA() = in(sa2_a, x: bitstring); 0.
let SA2ServerB() = in(sa2_b, x: bitstring); 0.

let VoterPrivacy(skV: skey, pkEA: pkey) =
    new r_enc: nonce;
    new r_com: nonce;
    new r_sa2: nonce;
    let v = diff[v0, v1] in
    let encVote = enc(v, pkEA, r_enc) in
    let com = commit(v, r_com) in
    let zk = zkp_make(v, r_com, com) in
    let sig = ring_sign(encVote, skV, pk(skV)) in
    let ki = key_image(skV) in
    let b = make_ballot(encVote, com, zk, sig, ki) in
    out(sa2_a, sa2_share_a(encVote, r_sa2));
    out(sa2_b, sa2_share_b(encVote, r_sa2));
    out(bb, b);
    0.

process
    new skEA: skey;
    let pkEA = pk(skEA) in
    out(ch, pkEA);
    (
        new skAlice: skey;
        new skBob: skey;
        out(ch, pk(skAlice));
        out(ch, pk(skBob));
        VoterPrivacy(skAlice, pkEA) |
        VoterPrivacy(skBob, pkEA) |
        SA2ServerA() |
        SA2ServerB()
    )
\end{verbatim}
\end{footnotesize}

\subsection*{ProVerif output (excerpt)}
\begin{footnotesize}
\begin{verbatim}
$ proverif privacy.pv
...
RESULT Observational equivalence is true.
--------------------------------------------------------------
Verification summary:

Observational equivalence is true.

--------------------------------------------------------------
\end{verbatim}
\end{footnotesize}

\textbf{Interpretation.} The ``\texttt{Observational equivalence is true}'' result means that for every possible adversary strategy, the trace of network observations is identical regardless of whether each voter cast \texttt{v0} or \texttt{v1}. Therefore an adversary cannot infer how any individual voter voted from network-level observations alone.

\section{Individual Verifiability: \texttt{verifiability.pv}}
\label{sec:appx-proverif-verifiability}

Individual verifiability is verified via injective correspondence: every recorded ballot on the bulletin board corresponds to exactly one authentic cast event by a registered voter. The query asserts \texttt{BallotRecorded(b)} $\Longrightarrow$ \texttt{inj-event(VoterCastBallot(pkV, b))}.

\subsection*{Source: \texttt{proverif/verifiability.pv}}
\begin{footnotesize}
\begin{verbatim}
(* CovertVote: Individual Verifiability *)

type skey.
type pkey.
type vote.
type nonce.
type commitment.
type zkprf.
type signature.
type keyimage.

free ch: channel.
free bb: channel.
free auth_bb: channel [private].
free sa2_a: channel [private].
free sa2_b: channel [private].

free v0: vote.

fun pk(skey): pkey.
fun enc(vote, pkey, nonce): bitstring.
fun commit(vote, nonce): commitment.
fun zkp_make(vote, nonce, commitment): zkprf.
fun ring_sign(bitstring, skey, pkey): signature.
fun key_image(skey): keyimage.
fun sa2_share_a(bitstring, nonce): bitstring.
fun sa2_share_b(bitstring, nonce): bitstring.
fun make_ballot(bitstring, commitment, zkprf, signature, keyimage): bitstring.

event VoterCastBallot(pkey, bitstring).
event BallotRecorded(bitstring).

query b: bitstring, pkV: pkey;
    event(BallotRecorded(b)) ==> inj-event(VoterCastBallot(pkV, b)).

let SA2ServerA() = in(sa2_a, x: bitstring); 0.
let SA2ServerB() = in(sa2_b, x: bitstring); 0.

let VoterVerifiable(skV: skey, v: vote, pkEA: pkey) =
    new r_enc: nonce;
    new r_com: nonce;
    new r_sa2: nonce;
    let encVote = enc(v, pkEA, r_enc) in
    let com = commit(v, r_com) in
    let zk = zkp_make(v, r_com, com) in
    let sig = ring_sign(encVote, skV, pk(skV)) in
    let ki = key_image(skV) in
    let b = make_ballot(encVote, com, zk, sig, ki) in
    event VoterCastBallot(pk(skV), b);
    out(sa2_a, sa2_share_a(encVote, r_sa2));
    out(sa2_b, sa2_share_b(encVote, r_sa2));
    out(auth_bb, b);
    0.

let BulletinBoard() =
    in(auth_bb, b: bitstring);
    event BallotRecorded(b);
    out(bb, b);
    0.

process
    new skEA: skey;
    let pkEA = pk(skEA) in
    out(ch, pkEA);
    (
        new skV: skey;
        VoterVerifiable(skV, v0, pkEA) |
        BulletinBoard() |
        SA2ServerA() |
        SA2ServerB()
    )
\end{verbatim}
\end{footnotesize}

\subsection*{ProVerif output (excerpt)}
\begin{footnotesize}
\begin{verbatim}
$ proverif verifiability.pv
...
Starting query inj-event(BallotRecorded(b_2))
              ==> inj-event(VoterCastBallot(pkV,b_2))
goal reachable: b-inj-event(VoterCastBallot(...))
              -> inj-event(BallotRecorded(...))
The hypothesis occurs strictly before the conclusion.

RESULT inj-event(BallotRecorded(b_2)) ==>
       inj-event(VoterCastBallot(pkV,b_2)) is true.
--------------------------------------------------------------
Verification summary:

Query inj-event(BallotRecorded(b_2)) ==>
      inj-event(VoterCastBallot(pkV,b_2)) is true.

--------------------------------------------------------------
\end{verbatim}
\end{footnotesize}

\textbf{Interpretation.} The injective correspondence guarantees that no recorded ballot is forged or duplicated: each \texttt{BallotRecorded} event on the bulletin board is matched to a unique prior \texttt{VoterCastBallot} event by an authenticated voter, providing individual verifiability.

\section{Voter Eligibility: \texttt{eligibility.pv}}
\label{sec:appx-proverif-eligibility}

Voter eligibility is verified via causal correspondence: every valid ballot cast was preceded by a legitimate registration event with the registration authority. The query asserts \texttt{ValidBallotCast(pkV)} $\Longrightarrow$ \texttt{VoterRegistered(pkV)}.

\subsection*{Source: \texttt{proverif/eligibility.pv}}
\begin{footnotesize}
\begin{verbatim}
(* CovertVote: Voter Eligibility *)
(* NO diff allowed -- query mode only *)

type skey.
type pkey.
type vote.
type credential.
type nonce.

free ch: channel.
free bb: channel.
free reg: channel [private].

free v0: vote.

fun pk(skey): pkey.
fun enc(vote, pkey, nonce): bitstring.
fun real_cred(skey): credential.
fun fake_cred(skey): credential.

event VoterRegistered(pkey).
event ValidBallotCast(pkey).

query pkV: pkey;
    event(ValidBallotCast(pkV)) ==> event(VoterRegistered(pkV)).

let RegistrationAuthority() =
    in(reg, voterPK: pkey);
    new skCred: skey;
    out(reg, (real_cred(skCred), fake_cred(skCred)));
    0.

let VoterEligible(skV: skey, v: vote, pkEA: pkey) =
    out(reg, pk(skV));
    in(reg, (realC: credential, fakeC: credential));
    event VoterRegistered(pk(skV));
    new r_enc: nonce;
    let encVote = enc(v, pkEA, r_enc) in
    event ValidBallotCast(pk(skV));
    out(bb, encVote);
    0.

process
    new skEA: skey;
    let pkEA = pk(skEA) in
    out(ch, pkEA);
    (
        new skV: skey;
        RegistrationAuthority() |
        VoterEligible(skV, v0, pkEA)
    )
\end{verbatim}
\end{footnotesize}

\subsection*{ProVerif output (excerpt)}
\begin{footnotesize}
\begin{verbatim}
$ proverif eligibility.pv
...
Starting query event(ValidBallotCast(pkV))
              ==> event(VoterRegistered(pkV))
goal reachable: b-event(VoterRegistered(pk(skV[])))
              && mess(reg[],voterPK_1)
              -> event(ValidBallotCast(pk(skV[])))

RESULT event(ValidBallotCast(pkV)) ==>
       event(VoterRegistered(pkV)) is true.
--------------------------------------------------------------
Verification summary:

Query event(ValidBallotCast(pkV)) ==>
      event(VoterRegistered(pkV)) is true.

--------------------------------------------------------------
\end{verbatim}
\end{footnotesize}

\textbf{Interpretation.} The causal correspondence guarantees that any ballot accepted as ``valid'' is necessarily preceded by a successful registration of the same voter public key with the registration authority. This rules out unauthorised vote injection by unregistered parties.

\section{Summary}
\label{sec:appx-proverif-summary}

Table~\ref{tab:proverif} of Chapter~\ref{ch:results} summarises the three results in tabular form: all three properties hold under the Dolev--Yao adversary model. Together they establish symbolic-model security against eavesdropping, replay, and man-in-the-middle attacks at the protocol level. Stronger guarantees in the computational model (game-based or universally composable proofs) are identified as future work in Chapter~\ref{ch:conclusion}.
